% ===============================================================
% Section VI  --  Results and Analysis
% Owner: Evaluation Lead (integrated)
% ===============================================================
\section{Results and Analysis}
\label{sec:results}

\subsection{Overall Results}
Table~\ref{tab:results-summary} reports average scores across all 15
questions. Retrieval relevance is not applicable to the baseline, which
does not retrieve.

\begin{table}[ht]
\centering
\caption{Average scores across 15 questions (scale 1--5). Generation was
run with the stand-in generator described in
Section~\ref{sec:evaluation}.}
\label{tab:results-summary}
\begin{tabular}{lcccc}
\toprule
System & Correct. & Ground. & Retrieval & Useful. \\
\midrule
Baseline           & 2.47 & N/A  & N/A  & 2.47 \\
RAG (scene)        & 2.92 & 2.00 & 5.00 & 2.92 \\
Enhanced RAG (utt.)& 2.80 & 1.87 & 5.00 & 2.80 \\
\bottomrule
\end{tabular}
\end{table}

Two findings stand out. First, \textbf{retrieval works}: both RAG systems
scored a perfect 5.00 on retrieval relevance, meaning the right scenes
and utterances were consistently surfaced for each question. Since
retrieval is independent of the generation model, this result transfers
directly to the deployed Gemma-based system. Second, \textbf{the
generation stage is the bottleneck}: grounding scores were low (2.00 and
1.87) even though the correct passages were retrieved, indicating that
the stand-in generator frequently failed to use the evidence it was
given. The RAG systems still edged out the baseline on correctness and
usefulness, confirming that retrieval helps even a weak generator, but
the gap is smaller than it would be with a capable model.

The utterance-level enhanced system did not clearly beat the scene-level
system; it scored slightly lower on grounding (1.87 vs 2.00). Individual
utterances, stripped of surrounding scene context, appear to give a weak
generator less to work with, not more, which is consistent with the
chunking trade-off anticipated in Section~\ref{sec:dataset}.

\subsection{Qualitative Examples}

\textbf{Q01 --- Who is Hamlet? (retrieval helps).} The baseline looped its
own prompt back into the output, producing a non-answer. The RAG system
retrieved Act~4 Scene~3 and Act~5 Scene~1, the correct scenes, and
produced a partial answer that at least reflected retrieved content,
earning a higher correctness score.

\textbf{Q02 --- Role of Lady Macbeth? (retrieval works, generation
fails).} The baseline hallucinated entirely. The RAG system retrieved the
murder scene (Act~2 Scene~2) correctly but the generator failed to use
it, repeating a nonsensical phrase. This is a clear instance of correct
retrieval undone by weak generation.

\textbf{Q13 --- Stylised generation.} All three systems struggled. The
stand-in generator is too small to produce convincing Early-Modern
register; retrieval supplied relevant context, but the generated style
was weak across systems.

\subsection{Failure Analysis}

\textbf{1. Generation model too weak.} The most consistent failure was
the generation step in the batch harness. The stand-in generator
frequently looped its prompt or produced repetitive, incoherent text even
though retrieval scored 5.00. This is the single largest driver of the
low correctness and grounding scores and is precisely the gap that the
deployed Gemma~3~4B model is intended to close.

\textbf{2. Empty or broken outputs on some scene-level questions.} Q03
(Montagues vs.\ Capulets) and Q12 (how Macbeth changes) produced no
scores for the scene-level system due to a generation failure rather than
a retrieval failure, as the same questions were answered by the
utterance-level system. \emph{Improvement:} validate generator output and
retry with a fallback prompt when an empty or degenerate response is
detected; the deployed Gemma model removes the underlying cause.

\textbf{3. Low grounding despite correct retrieval.} Both RAG systems
scored 5.00 on retrieval but around 2.00 on grounding: the right passages
were found but not incorporated. The generator effectively fell back on
its own (very limited) prior knowledge rather than the supplied evidence.
\emph{Improvement:} a stronger instruction-following generator (Gemma~3~4B)
plus a prompt that requires the answer to cite the supplied passages
directly closes this retrieval-to-generation gap.

\textbf{4. Enhanced RAG did not outperform scene-level RAG.} We expected
utterance-level retrieval to win by pinpointing exact lines; in practice
it scored marginally lower on grounding. Without surrounding scene
context, isolated utterances gave a weak generator less to anchor on.
\emph{Improvement:} attach each retrieved utterance's parent-scene summary
to the context block so the precision of utterance retrieval is combined
with the context of scene-level chunks.

\subsection{Summary}
The retrieval component is reliable across both granularities; the
limiting factor in the batch evaluation was the stand-in generator. The
highest-impact improvement, and the one that aligns the evaluation with
the demonstrable system, is to re-run the harness against Gemma~3~4B so
that the strong retrieval scores are matched by coherent, grounded
answers. The failure modes above are written so that this re-run updates
the numbers without changing the analysis structure.
